\chapterimage{chapter_head_2.pdf}
\chapter{Algorithm Analysis \& Algorithm Design}

\section{Theoretical Core: Big-O, Omega, and Theta}

\begin{definition}[Big-O Notation ($O$)]
Represents the asymptotic upper bound of an algorithm's runtime or space. It describes the worst-case scenario, guaranteeing the algorithm will take no more than a certain amount of time for sufficiently large inputs.
\end{definition}

\begin{definition}[Omega Notation ($\Omega$)]
Represents the asymptotic lower bound. It describes the best-case scenario, guaranteeing the algorithm will take at least a certain amount of time.
\end{definition}

\begin{definition}[Theta Notation ($\Theta$)]
Represents the asymptotically tight bound. It is used when an algorithm is bounded both above and below by the same function (i.e., both Big-O and Omega are the same).
\end{definition}

\newpage
\section{Common Time Complexities}

\begin{theorem}[Hierarchy of Time Complexities]
Common algorithm runtimes categorized from most efficient to least efficient:
\begin{itemize}
    \item \textbf{$O(1)$ - Constant Time:} Runtime is independent of the input size. Example: Array index access.
    \item \textbf{$O(\log n)$ - Logarithmic Time:} Runtime grows logarithmically. Example: Binary search on a sorted list.
    \item \textbf{$O(n)$ - Linear Time:} Runtime grows proportionally with the input size. Example: Linear search through an array.
    \item \textbf{$O(n \log n)$ - Linearithmic Time:} Typical of optimal comparison-based sorts. Example: Merge sort or the binary search intersection strategy.
    \item \textbf{$O(n^2)$ - Quadratic Time:} Runtime grows with the square of the input size. Example: Comparing all pairs in an array or a nested loop linear search intersection strategy.
\end{itemize}
\end{theorem}

\section{Analyzing Loops and Nested Loops}

\subsection{Single Loops}
The runtime correlates to the number of iterations times the work done inside the loop. For $n$ iterations doing constant work, the time is $O(n)$.

\subsection{Nested Loops}
Evaluated by multiplying the number of executions of the inner loop by the outer loop (when independent). For example, an outer loop running $n$ times and an inner loop running $m$ times results in $O(n \times m)$ time.

\subsection{Algorithmic Tracing: Sorted Array Matching}
When counting common elements in two sorted arrays of sizes $n$ and $m$, different strategies yield vastly different runtimes:

\begin{example}[Strategy 1: Linear Search]
An outer loop over $n$ elements and an inner loop over $m$ elements yields $O(n \times m)$.
\end{example}

\begin{example}[Strategy 2: Binary Search]
A loop over $n$ elements performing a binary search in $m$ elements yields $O(n \log m)$.
\end{example}

\begin{example}[Strategy 3: Two Pointers]
Iterating both arrays simultaneously using two index variables yields an efficient $O(n + m)$.
\end{example}

\newpage
\section{Code Implementation}

The Two-Pointer Strategy is the most efficient way to find the intersection of two sorted arrays.

\begin{lstlisting}[language=C, caption=Two-Pointer Strategy for Sorted Array Intersection]
#include <stdio.h>

// Returns the number of common elements between two sorted arrays
int count_common(int arr1[], int n, int arr2[], int m) {
    int i = 0, j = 0;
    int count = 0;

    // Advance the index of the smaller value, count when equal
    while (i < n && j < m) {
        if (arr1[i] == arr2[j]) {
            count++;
            i++;
            j++;
        } else if (arr1[i] < arr2[j]) {
            i++;
        } else {
            j++;
        }
    }
    return count;
}
\end{lstlisting}

\section{Skills \& Pitfalls}

\subsection{Must-Know Skills}
\begin{itemize}
    \item \textbf{Empirical Runtime Calculation:} Using experimental elapsed times to find the theoretical constant $C$. Given $T(n) = C \cdot f(n)$, we can calculate $C = T(n) / f(n)$. This enables the prediction of runtimes for arbitrary input sizes.
    \item \textbf{Solving Constraints:} Re-arranging empirical equations to find the maximum allowed input size $N$ for a given constraint $T$.
\end{itemize}

\subsection{Common Mistakes (Pitfalls)}

\begin{remark}
\textbf{Ignoring Constant Factors:} While $O(N)$ ignores constants theoretically, in practice (empirical runtime), the constant factor dictates performance for realistic input sizes.
\end{remark}

\begin{remark}
\textbf{Misidentifying the Bottleneck:} Failing to recognize which loop or operation dominates the overall runtime (e.g., misidentifying the bottleneck as a single $O(N)$ loop when a hidden $O(N^2)$ operation is executed).
\end{remark}

\begin{remark}
\textbf{Assuming Equal Lengths:} Assuming inputs (like two arrays) always have the same length $N$, ignoring edge case runtimes like $O(M \times N)$ or $O(M \log N)$.
\end{remark}
